\title{\textbf{Autonomous Prime-Valuation Dynamics and\\ Symbolic Complexity in the UNNS Substrate}\\[6pt]
\large Prime-fire reduction, an irrational rotation and Sturmian itineraries, prime-support renewal, finite-guard complexity and arithmetic counters}
\author{UNNS Research Collective}
\date{Version 1.2 \,--\, 11 November 2026}

\begin{document}
\maketitle

\begin{abstract}
\noindent
We study the autonomous arithmetic dynamics generated by the forward transitions $x\mapsto x(N+1)^2/N$, $N\mid x$, $N\ge2$, of the UNNS Substrate.
In the coordinates $v_p(x)$ the transition is the exact integer translation $v_p\mapsto v_p+2v_p(N+1)-v_p(N)$.
For every selector that fires primes we prove that primes $\ge5$ are fired only finitely often, so that the eventual dynamics lives on the two valuations $(\alpha,\beta)=(v_2,v_3)$ with translations $(-1,+2)$ and $(+4,-1)$.
In this reduced system the rule that fires the prime with the largest prime-power part (\textsf{MAXPP}) is an irrational rotation of a circle of length $\ln 864$ by the angle $\ln 48$, coded by an arc of length $\ln 18$.
Consequently, for \emph{every} initial integer $x_0\ge2$ the selection itinerary of \textsf{MAXPP} is not eventually periodic and its eventual factor complexity is exactly $p(k)=k+1$, so it is a Sturmian word of entropy $0$.
The frequency of prime $2$ is $\theta=\ln48/\ln864\approx0.572530$ with discrepancy below one, the valuations are Beatty-type sequences, and $\log_2x_n$ grows at the rate $2.2747033$ bits per step.
Allowing composite branch parameters, we prove that prime-support renewal occurs if and only if infinitely many distinct branch values are used, and that for selectors with a finite set of branch values, affine valuation guards, thresholds, congruences and finite control the branch/control itinerary has polynomial factor complexity and zero entropy.
We also give an autonomous state-only selector whose branch values form the counter $m_0,m_0+1,\dots$ and for which an exponent-profile observable has entropy at least $2\ln2/\pi^2$; this separates the complexity \emph{generated} by selection from the arithmetic complexity \emph{revealed} by an observable.
All computations are reproduced here by newly written code; the earlier releases are inherited and not re-run.
\end{abstract}

\tableofcontents
\newpage

%=====================================================================
\section{Introduction}\label{sec:intro}
%=====================================================================
\subsection{The setting}
The UNNS Substrate contains an operation that nests a number $N$ into a multiple $x$ of it and replaces $x$ by
\[
T_N(x)\;=\;x\,\frac{(N+1)^2}{N},\qquad N\mid x,\ N\ge2 .
\]
Iterating it requires a choice of $N$ at each step.
A \emph{selector} makes this choice; it is \emph{autonomous} if the choice is a function of the current integer alone, with no clock, no external input and no memory.
The sequence of choices is the \emph{itinerary}, and the questions of this paper are about the itinerary: is it eventually periodic, how many distinct words of each length does it contain, and what mechanism produces the answer.

The natural coordinates for these questions are the prime valuations $v_p(x)$.
In these coordinates a transition is an integer translation (Theorem~\ref{thm:valuation}), and the whole theory is the study of a system of integer-coordinate translations that are admissible only where $v_p(x)\ge v_p(N)$ and that are chosen by an autonomous rule.

\subsection{Results}
\begin{theorem}[Summary]\label{thm:summary}
\begin{enumerate}[label=(\Alph*),itemsep=2pt]
\item \textbf{Prime-fire reduction} (Theorem~\ref{thm:A}; \cite{u4gen}). For every selector that fires primes and every $x_0\ge2$, primes $\ge5$ are fired only finitely often. The eventual dynamics is the two-valuation system $(\alpha,\beta)=(v_2,v_3)$ with the translations $2:(\alpha,\beta)\mapsto(\alpha-1,\beta+2)$ and $3:(\alpha,\beta)\mapsto(\alpha+4,\beta-1)$.
\item \textbf{Autonomous Sturmian generation} (Theorem~\ref{thm:B}; \cite{u4gen}). The rule \textsf{MAXPP} (fire the prime with the largest prime-power part) generates from every $x_0\ge2$ an itinerary that is not eventually periodic and whose eventual factor complexity is exactly $p(k)=k+1$; its entropy is $0$.
\item \textbf{Exact frequencies and growth} (Theorem~\ref{thm:C}; \cite{u4genc}). The frequency of the prime $2$ in the \textsf{MAXPP} itinerary is $\theta=\ln48/\ln864\approx0.572530$, the number of $2$'s among the first $n$ letters of the tail differs from $n\theta$ by less than $1$, and $\log_2x_n$ grows at the rate $\theta\log_2\tfrac92+(1-\theta)\log_2\tfrac{16}3\approx2.2747033$ bits per step.
\item \textbf{Prime-support renewal} (Theorem~\ref{thm:D}; \cite{u4genc}). Along a forward trajectory the total prime support is finite if and only if only finitely many distinct branch values $N$ occur.
\item \textbf{Symbolic complexity} (Theorem~\ref{thm:E}; \cite{u4genc}, scope as in Definition~\ref{def:finite-guard}). For selectors with finitely many branch values, affine valuation guards, thresholds, congruences and finite control, the branch/control itinerary has polynomial factor complexity and zero entropy.
\item \textbf{Arithmetic counters} (Theorems~\ref{thm:ordS}, \ref{thm:ordE}; \cite{u4order}). There is an autonomous state-only selector whose branch values are $N_j=m_0+j$, $x_j=(m_0+j)(m_0+j)!/m_0!$, and an exponent-profile observable of its itinerary with entropy $\ge2\ln2/\pi^2$.
\end{enumerate}
\end{theorem}
Items (A)--(C) are exact statements about integer dynamics; (D) and (E) describe what changes when composite branch values are allowed; (F) shows how an observable can acquire positive entropy from the arithmetic of a counter rather than from the selection dynamics.
Section~\ref{sec:signif} states their combined meaning.

\subsection{Evidence levels}\label{ssec:evidence}
Three kinds of statement are distinguished throughout.
\begin{itemize}[itemsep=1pt]
\item \textbf{Proved}: a complete argument is given in the text or in an appendix. Classical results used without proof are named: Størmer's theorem on consecutive smooth integers (Theorem~\ref{thm:stormer}) and Euler's product for $\zeta(2)$.
\item \textbf{Computed}: an exact computation with integers or with multiprecision arithmetic reproduced by the code of this release (Appendix~\ref{app:checks}). Computations illustrate and test proved statements; none is used to prove an asymptotic claim.
\item \textbf{Inherited}: a result of an earlier release that is cited with its archive hash and is not re-run here (Appendix~\ref{app:prov}, which gives the source of each theorem family; the sources are \cite{u4gen,u4genc,u4order}).
\end{itemize}

\subsection{Organisation}
Section~\ref{sec:grammar} fixes the transition grammar and the valuation law. Section~\ref{sec:fire} proves the prime-fire reduction and treats the periodic selector regimes. Sections~\ref{sec:maxpp}--\ref{sec:freq} are the centre of the paper: the \textsf{MAXPP} selector, its description as an irrational rotation, the Sturmian complexity, and the exact frequencies and growth. Section~\ref{sec:composite} treats composite branches and prime-support renewal, Section~\ref{sec:complexity} the finite-guard complexity theorem, and Section~\ref{sec:order} arithmetic counters and exponent-profile observables. Section~\ref{sec:signif} states the mathematical significance. The appendices contain the classical lemmas, the certificates and the provenance of the sources.

%=====================================================================
\section{Transition grammar and prime-valuation coordinates}\label{sec:grammar}
%=====================================================================
\subsection{The forward transition}
For integers $x\ge2$ and $N\ge2$ with $N\mid x$ put $T_N(x)=x(N+1)^2/N$. In the UNNS arithmetic graph, whose edges are $\{kN,k(N+1)^2\}$, this is the forward edge from $kN$ to $k(N+1)^2$ with $k=x/N$; only its arithmetic is used below.
\begin{proposition}[Legality and growth]\label{prop:growth}
If $N\mid x$ then $T_N(x)$ is an integer and
\[
\frac{T_N(x)}{x}=N+2+\frac1N\ \ge\ \frac92 .
\]
In particular $T_N(x)>4x$, so a forward orbit is strictly increasing, never returns to a vertex, and $x_n\ge 2\cdot 4.5^n$.
\end{proposition}
\begin{proof} $x/N$ is an integer, hence so is $(x/N)(N+1)^2$. The ratio is $(N+1)^2/N=N+2+1/N$, which is increasing in $N$ and equals $9/2$ at $N=2$. \end{proof}

\subsection{The valuation law}
For a prime $p$ let $v_p(x)$ be the exponent of $p$ in $x$, and write $\supp(x)=\{p:v_p(x)>0\}$.
\begin{theorem}[Valuation law]\label{thm:valuation}
For every legal transition and every prime $p$,
\[
\boxed{\;v_p\bigl(T_N(x)\bigr)=v_p(x)+2\,v_p(N+1)-v_p(N).\;}
\]
The transition is legal at $x$ if and only if $v_p(x)\ge v_p(N)$ for all primes $p$.
\end{theorem}
\begin{proof} $T_N(x)=x\,(N+1)^2\,N^{-1}$ and $v_p$ is additive on products and negated on inverses. Legality $N\mid x$ is the statement $v_p(x)\ge v_p(N)$ for all $p$. \end{proof}
Since $\gcd(N,N+1)=1$, for each prime $p$ at most one of $v_p(N),v_p(N+1)$ is non-zero: a transition lowers the exponents of the primes of $N$ and raises those of $N+1$ by twice their exponent. Writing $x$ as the exponent vector $(v_p(x))_p$, every transition is the translation of this vector by $\tau_N=\bigl(2v_p(N+1)-v_p(N)\bigr)_p$, applied where it is admissible. The two primes that matter most are $2$ and $3$:
\begin{center}
\begin{tabular}{@{}lll@{}}\toprule
$N$ & translation of $(v_2,v_3)$ & other effect\\\midrule
$2$ & $(-1,+2)$ & none\\
$3$ & $(+4,-1)$ & none\\
$p\ge5$ prime & $(2v_2(p{+}1),\,2v_3(p{+}1))$ & $v_p-1$; $+2v_q(p{+}1)$ for $q\mid p{+}1$\\\bottomrule
\end{tabular}
\end{center}

\subsection{Selectors, orbits and itineraries}
\begin{definition}[Selector, orbit, itinerary]\label{def:selector}
A \emph{selector} chooses, at each step $n$, a branch value $N_n\ge2$ with $N_n\mid x_n$, and $x_{n+1}=T_{N_n}(x_n)$. The sequence $(N_n)_{n\ge0}$ is the \emph{itinerary} (or branch itinerary). A selector is \emph{autonomous} (or \emph{state-only}) if $N_n=\sigma(x_n)$ for a fixed function $\sigma$ of the integer alone. A step with $N_n$ prime is a \emph{prime-fire} step, and a selector all of whose steps are prime-fire is a \emph{prime-fire selector}; it may be of any kind, including one that consults the history.
\end{definition}
By Proposition~\ref{prop:growth} an orbit never repeats a vertex.

\subsection{Factor complexity and entropy}
For an infinite word $w=w_0w_1w_2\cdots$ over an alphabet and $k\ge1$, let $p_w(k)$ be the number of distinct factors $w_iw_{i+1}\cdots w_{i+k-1}$, $i\ge0$.
\begin{lemma}[Entropy]\label{lem:entropy}
$p_w(j+k)\le p_w(j)\,p_w(k)$. Hence $\ent(w)=\lim_{k\to\infty}\tfrac1k\ln p_w(k)$ exists and equals $\inf_k\tfrac1k\ln p_w(k)$. If $p_w(k)$ is bounded by a polynomial in $k$ then $\ent(w)=0$.
\end{lemma}
\begin{proof} A factor of length $j+k$ is the concatenation of a factor of length $j$ and one of length $k$, which gives the inequality. Then $a_k=\ln p_w(k)$ is subadditive and Fekete's lemma gives the limit and the infimum. Polynomial growth gives $a_k=O(\ln k)$. \end{proof}
The word $w$ is \emph{eventually periodic} if $w_{i+q}=w_i$ for some $q\ge1$ and all $i\ge i_0$.
